\section{Actual exponent realizations beyond the parity obstruction}
\label{lp-section}

For a primitive positive seed $(a,b)$ put
\[
 M_0=a^2+ab+b^2,\quad U=ab,\quad
 M_1=U^2+UM_0+M_0^2=F(a,b).
\]
The actual Eisenstein integers are $z_0=a+b\zeta$ and
$z_1=U+M_0\zeta$, where $\zeta^2-\zeta+1=0$.
Both are primitive, since $\gcd(ab,M_0)=1$.
We distinguish direct norm-residue tests, local power solutions, and
actual simultaneous extractions. None of the constructions below
supplies the small-lambda two-step gate.

\subsection{A precise boundary of direct fixed-modulus tests}

\begin{theorem}[Simultaneous direct power residues]\label{lp-residues}
Fix $m,n\ge2$ and finitely many positive integer moduli. There are
infinitely many primitive positive seeds, with $a/(a+b)\to1/2$ and
$a+b\to\infty$, for which $M_0$ and $M_1$ are respectively $m$-th
and $n$-th power residues modulo every specified modulus.
\end{theorem}
\begin{proof}
Let $L$ be the least common multiple of the moduli and take
$(a,b)=(Lk,Lk+1)$, $k\ge1$. These integers are coprime. Their
residues are $(0,1)$ modulo $L$, so $M_0\equiv M_1\equiv1\pmod L$.
The integer one is both specified power residues. The limiting height
and balance assertions follow directly.
\end{proof}
This excludes a universal obstruction consisting solely of finitely many
fixed power-residue tests on the original norm integers. It does not
exclude every argument using local information: the successful QG3
descent uses auxiliary covering classes and additional global input.
The seeds here are not claimed to have power norms. Indeed the theorem
also applies to the globally excluded square/square positive branch.

\begin{theorem}[Primitive local power solutions]\label{lp-local}
For every $m,n\ge2$ and every prime $p$, there are $a,b\in\mathbb Z_p$
with $ab(a+b)\ne0$, at least one coordinate a unit, such that $M_0$
is an $m$-th power and $M_1$ an $n$-th power in $\mathbb Z_p$.
The same system has a positive real solution.
\end{theorem}
\begin{proof}
Put $s=2\max\{v_p(m),v_p(n)\}+1$, $a=p^s$, and $b=a+1$.
For consecutive seeds, direct expansion gives
\begin{align}
 N_c(a)&=M_0=1+3a+3a^2,\label{lp-norm-polynomial}\\
 F_c(a)&=M_1=1+7a+20a^2+26a^3+13a^4.\label{lp-quartic-polynomial}
\end{align}
Both differences from one have valuation at least $s$. Applying strong
Hensel lifting at one to $T^m-M_0$ and $T^n-M_1$ is valid because
the error valuation exceeds twice the derivative valuation. It gives
unit roots in $\mathbb Z_p$. The chosen integers have the stated
nonzero and primitive properties. Over the reals, take the positive
roots of the norms of any positive seed.
\end{proof}
Different primes may use different seeds and roots. Local solubility
therefore does not supply a common rational or integer power solution.

\subsection{Prescribed actual depths with nonsquare residuals}

\begin{theorem}[Arbitrary two-step extraction exponents]\label{lp-actual}
For every $g_0,g_1\ge2$, there are infinitely many primitive positive
consecutive seeds, balanced in the limit, with actual representations
\[
 z_0=v_0(1+2\zeta)^{g_0},\qquad
 z_1=v_1(2+7\zeta)^{g_1}.
\]
Writing $V_i=N(v_i)$, the residual norms are positive nonsquares and
\begin{gather*}
 M_0=V_0\,7^{g_0},\qquad M_1=V_1\,67^{g_1},\\
 v_7(V_0)=v_{67}(V_1)=0,\qquad
 v_{13}(V_0)=v_{31}(V_1)=1.
\end{gather*}
In particular both exponents may be arbitrary odd numbers, or arbitrary
numbers congruent to two modulo four.
\end{theorem}
\begin{proof}
At $a=1$ the consecutive polynomials satisfy
\[
 N_c(1)=7,\quad N_c'(1)=9,\qquad
 F_c(1)=67,\quad F_c'(1)=177\equiv43\pmod{67}.
\]
Each is a simple root at its indicated prime. A simple root modulo
$p$ lifts uniquely to each modulus $p^e$. Of the $p$ extensions of
a root modulo $p^e$ to modulus $p^{e+1}$, exactly one remains a root;
choosing any other extension gives exact valuation $e$.
Thus choose a class modulo $7^{g_0+1}$ with
$v_7(N_c(a))=g_0$, and a class modulo $67^{g_1+1}$ with
$v_{67}(F_c(a))=g_1$. Both reduce to $a=1$ at their base prime.

Also impose $a\equiv5\pmod{13^2}$ and $a\equiv14\pmod{31^2}$.
The exact identities
\[
 N_c(5)=91=7\cdot13,\qquad
 F_c(14)=574771=31\cdot18541,\quad18541\equiv3\pmod{31}
\]
give residual-prime depth one. The four moduli are pairwise coprime,
so CRT provides an infinite progression of positive $a$ satisfying
all four exact valuation requirements. Set $b=a+1$.

The two indicated Eisenstein primes have norms seven and sixty-seven.
Modulo seven, $z_0$ reduces to $1+2\zeta$; modulo sixty-seven,
$z_1$ reduces to $2+7\zeta$. Thus each prime occurs with the stated
orientation. Primitivity prevents its conjugate from also dividing
the same integer. Unique factorization and the exact norm valuations
give the displayed powers of these fixed prime elements. Dividing
produces actual primitive Eisenstein quotients $v_i$, without any
assumption about the factorization of their remaining norms.

For consecutive seeds, $M_0\equiv1\pmod3$, so its ramified exponent
is zero; $M_1$ is also prime to three. Division by $7^{g_0}$ and
$67^{g_1}$ leaves the exact depths at thirteen and thirty-one, proving
that the residual norms are nonsquares. The exact root valuations
give their asserted coprimality with the displayed root norms.
Height and balance follow along the positive CRT progression.
\end{proof}

This is a complete integer realization beyond the square-residual parity
obstruction. It does not show that the residuals are analytically small.
At fixed $g_0,g_1$ their logarithms tend to infinity along the progression,
so both displayed $\lambda_i=\max\{1,\log V_i\}/g_i$ diverge.
For varying exponents, the earlier fixed-first-root obstruction still
applies whenever its bounded-first-lambda premise holds.

\begin{corollary}[Residual support escape in the realized family]
\label{lp-support-escape}
Along each progression in Theorem~\ref{lp-actual}, at fixed $g_0,g_1$,
the support of $V_1$ eventually fails to lie in each fixed finite prime
set. No effective rate of prime growth is asserted.
\end{corollary}
\begin{proof}
If $\operatorname{supp}(V_1)\subseteq S$, then
$\operatorname{supp}(M_1)\subseteq S\cup\{67\}$. The finite-total-
support theorem, Corollary~\ref{pl-total-support}, permits only finitely
many such primitive seeds. Their heights are eventually left by the
infinite positive progression.
\end{proof}
This last corollary uses the separate global finite-descent and Faltings
input. The residue and realization constructions themselves do not.

An exact integer replay verifies actual division and reconstruction in
the Eisenstein ring for twenty-one members with varied exponent pairs.
A separate scan of $76{,}116$ primitive pairs with $1\le a\le b\le500$
finds ten first-norm cubes and no second-norm cubes. The null scan is
only finite diagnostic evidence and proves no global nonexistence.
No statement in this section establishes high ABC quality, the required
two-step radical compression, or a proof or disproof of ABC.
